% =============================================================================
% Lecture 3: ML Lifecycle
% Introduction to Machine Learning Systems (Volume I)
% =============================================================================
% The lifecycle as a loop, not a pipeline. A decision made at any stage,
% above all at problem definition, constrains every other stage (the
% constraint propagation principle). Hook: Google's diabetic retinopathy
% screener, accurate in the lab, struggling in Thai clinics. The chapter's
% DR scenario carries the whole lecture; each of the six stages gets one
% picture and one failure.
%
% Build: make -C slides build-03_ml_lifecycle   (figures: images/make_figs.py)
% Plan:  lecture_map.yml in this folder.
% =============================================================================

\documentclass[aspectratio=169, 12pt]{beamer}
\usepackage{../../assets/beamerthememlsys}
\mlsyssetup{
  volume={Volume I},
  chapter={Chapter 3},
  chaptertitle={ML Lifecycle},
  logo={../../assets/img/logo-mlsysbook.png},
  instlogo={../../assets/img/logo-harvard.png},
}
\usepackage[T1]{fontenc}
\usepackage[scaled=0.9]{helvet}
\usepackage{courier}
\renewcommand{\familydefault}{\sfdefault}
\usepackage{booktabs,amsmath}
\usetikzlibrary{arrows.meta,positioning}
\graphicspath{{images/}}
\setcounter{mlsystotalsections}{5}

\title{ML Lifecycle}
\author{Vijay Janapa Reddi}
\institute{Harvard University}
\date{}

% A large answer choice for predict and concept-question slides
\newcommand{\choice}[2]{\textbf{\textcolor{mlsysaccent}{#1}}\enspace #2}
% A highlighted (revealed) answer choice
\newcommand{\rchoice}[2]{\colorbox{mlsysaccent!15}{\choice{#1}{#2}}}
% An artifact chip for the diagnose slide: kept (gray) or redone (red) on step 2
\newcommand{\chip}[2]{%
  \alt<2->{\colorbox{#2}{\strut\textcolor{white}{\textbf{#1}}}}%
          {\colorbox{vlightgray}{\strut\textcolor{darktext}{\textbf{#1}}}}}

\begin{document}

% =============================================================================
% TITLE
% =============================================================================
\mlsysnote{[1 min] Read the thesis aloud once. Lecture 2 asked where a model can run. Today asks when a team must know that, and what happens when it learns too late. One system carries the whole lecture: screening for diabetic retinopathy, an eye disease caused by diabetes, in rural clinics. The notes total about 70 minutes.}
\mlsystitle{ML Lifecycle}{An ML system is built around a loop of six stages, and a decision made at any stage, above all at problem definition, constrains every other stage.}{cover_ml_lifecycle.png}

% =============================================================================
% HOOK
% =============================================================================
\begin{frame}{A retina screener matched eye specialists in the lab}
  \begin{columns}[c]
    \begin{column}{0.50\textwidth}
      \centering\safeimg[width=\textwidth,height=5.0cm,keepaspectratio]{fundus-camera-topcon.jpg}
    \end{column}
    \begin{column}{0.47\textwidth}
      \normalsize
      \textbf{In the lab} (2016)\\[0.1cm]
      128,000 graded photos\\
      Sensitivity \textbf{97.5\%}, specificity \textbf{93.4\%}\\[0.5cm]
      \uncover<2->{\textbf{In the clinic}\\[0.1cm]
      Deployed in Thailand; observed in 11 clinics over eight months}
    \end{column}
  \end{columns}
  \mlsyscredit{Photo: fundus camera, Ralf Roletschek (FAL), Wikimedia Commons. Lab: Gulshan et al., JAMA 2016. Clinics: Beede et al., CHI 2020; Google, ``Healthcare AI systems that put people at the center,'' 2020.}
  \note{[2 min] The photo is a fundus camera: it photographs the back of the eye, the retina, where diabetes damages small blood vessels. Diabetic retinopathy is a major cause of preventable vision loss, and screening needs eye specialists that rural regions lack. Step 1: Google trained a model on 128,000 retinal photos graded by ophthalmologists. Sensitivity is the share of sick patients the model flags; specificity is the share of healthy patients it clears. At one threshold it reached 97.5 and 93.4 percent, at the level of specialists. Step 2: Google Health then deployed it in clinics in Thailand, and a team spent eight months visiting 11 of them to watch nurses use it. ENGAGE: the next slide asks what happened. FLEX: [CORE].}
\end{frame}

\begin{frame}{Predict: what went wrong in the clinics?}
  \mlsysactivity{Predict}{2.5 min}
  \centering
  {\normalsize Same model, same task. The lab result did not carry over.}\\[0.45cm]
  \begin{minipage}{0.72\textwidth}
    \large
    \choice{A}{The model was less accurate on Thai patients}\\[0.15cm]
    \choice{B}{Many photos could not be graded}\\[0.15cm]
    \choice{C}{Uploads stalled on slow clinic internet}\\[0.15cm]
    \alt<2->{\rchoice{D}{B and C}}{\choice{D}{B and C}}
  \end{minipage}\\[0.4cm]
  \uncover<2->{\normalsize\textbf{393 of 1,838} photos (21\%) failed the image-quality check.}
  \mlsyscredit{Beede et al., CHI 2020, as cited in ML Lifecycle, \S3.1.2.}
  \note{[2.5 min] ENGAGE: vote with fingers, one to four; ask one student for A and one for D to explain. Step 2 reveals D. In the study, 393 of 1,838 submitted photos, 21 percent, were rejected as ungradable, often because clinic lighting, camera maintenance, or whether pupils were dilated did not match the system's assumptions. Rejected photos meant retakes, which added work for nurses, and slow clinic internet held up uploads to the remote grading service. WARN: most students pick A, because they picture the model as the thing that fails. The model was never the problem the study found. Every failure sat at a different stage of the system around it. That is today's subject. Do not assign the failures to stages yet; the first segment ends with that exercise. FLEX: [CORE].}
\end{frame}

\begin{frame}{Recall: in what order do we screen a deployment tier?}
  \mlsysactivity{Retrieval}{1 min}
  \centering
  {\Large\bfseries Close your notes.}\\[0.3cm]
  {\normalsize Write the four screening gates from Lecture 2, in order.}\\[0.6cm]
  \uncover<2->{\Large Privacy $\rightarrow$ latency $\rightarrow$ compute $\rightarrow$ cost}
  \mlsyscredit{System Paradigms, \S2.9.2.}
  \note{[1 min] Thirty seconds of silent writing, then step 2. A legal ban on moving data rules out options before anything is measured; a network round trip longer than the deadline rules out the cloud; the model must fit and run on what remains; only then is cost compared. LINK: today the clinic's in-visit deadline and its fragile uplink settle the tier at problem definition, and cost only confirms it at the end of the lecture. FLEX: [CORE].}
\end{frame}

\begin{frame}{Why see the whole map before walking any single path?}
  \normalsize By the end of today, you will be able to:\\[0.25cm]
  \begin{itemize}\setlength{\itemsep}{0.25cm}
    \item Trace how a decision at one stage constrains every other stage.
    \item Estimate the rework and the iteration cost of a late constraint.
    \item Turn a request into numeric targets: sensitivity, specificity, memory, latency.
    \item Decide where a model runs: feasibility first, then cost.
  \end{itemize}
  \note{[1 min] The title is the driving question, the one the chapter opens with. Read the four outcomes. Say that the whole lecture follows one system, the retinopathy screener, so students can watch a single decision travel through every stage. FLEX: [CORE].}
\end{frame}

% =============================================================================
% SECTION: THE LOOP
% =============================================================================
\mlsysnote{[0.5 min] Start with the map itself: six stages and what passes between them.}
\mlsysdivider{Why is the ML lifecycle a loop, not a line?}{ML Lifecycle, \S3.1}
\section{The Loop}

\begin{frame}{Six stages, each handing an artifact to the next}
  \centering
  \mlsyssteps[width=0.92\textwidth,height=5.4cm,keepaspectratio]{loop_step}{3}
  \note{[3 min] This diagram returns all lecture. Step 1: the six stages in the order a project first meets them. Define the problem, collect and prepare data, develop and train the model, evaluate and validate it, deploy and integrate it, monitor and maintain it. Step 2: what each stage hands the next. Problem definition produces measurable targets and limits; data collection produces a versioned dataset; training produces weights; validation produces a release manifest, the record of evidence that a specific model version passed; deployment produces a running service. Step 3: the dashed green arrow. Monitoring sends evidence back, about drifting inputs, missed cases, and new limits, and it can re-enter at any earlier stage. ENGAGE: ``Which arrow would a traditional software project draw thinly or not at all?'' (The return arrow.) WARN: students read the numbers as a schedule to finish once. The return arrow is what makes it a lifecycle: deployment starts a phase rather than ending a project. LINK: the operations loop that runs this return path is Lecture 14. FLEX: [CORE].}
\end{frame}

\begin{frame}{ML behavior can change when only the data changes}
  \centering\small
  \renewcommand{\arraystretch}{1.5}
  \begin{tabular}{@{}P{0.20\textwidth}P{0.34\textwidth}P{0.38\textwidth}@{}}
    \toprule
    & \textbf{Traditional software} & \textbf{ML system} \\
    \midrule
    Behavior comes from & code and configuration & code \textbf{and data} \\
    After deployment & fixed until code changes & \textcolor{RedLine}{\textbf{can degrade with code unchanged}} \\
    Production feedback & changes code and requirements & also becomes \textbf{training data} \\
    \bottomrule
  \end{tabular}
  \mlsyscredit{After ML Lifecycle, table 1.}
  \note{[1.5 min] Why a loop rather than a line. Both kinds of software iterate; the difference is where behavior comes from. A rule-based program changes behavior when someone edits code, configuration, or a dependency. An ML system learns its behavior from data, so it can change when the world changes while every line of code stays fixed. Example: a clinic replaces its camera, the new photos have a different color profile, and sensitivity falls with no code change and no error raised. Only fresh labeled outcomes reveal that, which is why monitoring must feed back. LINK: Lecture 1 showed silent degradation; this is its workflow consequence. FLEX: [OPTIONAL]; if short, say the one sentence ``ML behavior can change when only the data changes'' on the loop slide and skip this one.}
\end{frame}

\begin{frame}{The day-153 failure: a limit found after the work was done}
  \centering
  \mlsyssteps[width=0.98\textwidth,height=3.8cm,keepaspectratio]{day153_step}{3}\\[0.25cm]
  \uncover<3->{\normalsize Excellent accuracy. Excellent ML skills. A \textbf{workflow} failure.}
  \mlsyscredit{ML Lifecycle, \S3.1.2 (a scenario built on the deployment case).}
  \note{[3 min] The chapter's scenario: a team runs the stages as a line. Step 1: day 1, the request is one sentence, ``build a diabetic retinopathy screening model for rural clinics.'' Day 90, 92 percent sensitivity on the validation set; day 120, 94 percent after a month of tuning. Things look good. Step 2: day 150, the model is handed to deployment engineers. Day 151, they report it needs 4 GB of memory. Day 152, someone finally checks the device: a low-power edge computer in each clinic whose 4 to 8 GB of memory is shared with the operating system, the camera pipeline, and the clinic software, and the clinic's IT plan reserves about 500 MB for the model. Step 3: day 153, five months of work is discarded. ENGAGE: ``On what day should the 500 MB number have been written down?'' (Day 1.) WARN: students call this an accuracy problem or a compression problem. The model was good and the team was skilled; the limit simply arrived after the decisions that depended on it. FLEX: [CORE].}
\end{frame}

\begin{frame}{The model needed eight times its memory reservation}
  \centering
  \safeimg[width=0.92\textwidth,height=4.4cm,keepaspectratio]{memory_fit.pdf}\\[0.2cm]
  {\normalsize Compact models built for budgets like this exist. Nobody chose one.}
  \mlsyscredit{ML Lifecycle, \S3.1.2 and \S3.4; MobileNetV2: Sandler et al., CVPR 2018.}
  \note{[1 min] Log scale. Red: the day-151 model at 4 GB, eight times the reservation. Gray: about 500 MB reserved for the model. Orange: MobileNetV2, a network designed for phones and the book's mobile lighthouse model, has about 14 MB of weights in 32-bit floating point (FP32, four bytes per weight) and about 600 million operations per inference. WARN: weights are not the whole budget; activations (the intermediate values the network computes), input buffers, and the runtime also need memory, and the capture resolution needed to see small lesions raises the cost. The point is that the right family of models was available on day 1 if the limit had been known. LINK: Lecture 6 builds MobileNet; Lecture 10 shrinks models to fit. FLEX: [CORE].}
\end{frame}

\begin{frame}{Diagnose: which stage owned each failure?}
  \mlsysactivity{Diagnose}{2.5 min}
  \begin{columns}[c]
    \begin{column}{0.36\textwidth}
      \centering\safeimg[width=\textwidth,height=3.2cm,keepaspectratio]{loop_step3.pdf}
    \end{column}
    \begin{column}{0.62\textwidth}
      \normalsize
      \renewcommand{\arraystretch}{1.35}
      \begin{tabular}{@{}P{0.52\textwidth}P{0.40\textwidth}@{}}
        \toprule
        \textbf{Clinic symptom} & \textbf{Stage} \\
        \midrule
        21\% of photos ungradable & \leavevmode\uncover<2->{\textcolor{GreenLine}{\textbf{2 Data}}, \textbf{4 Validate}} \\
        Retakes add nurse work & \leavevmode\uncover<3->{\textbf{1 Define}, \textbf{4 Validate}} \\
        Uploads stall on slow links & \leavevmode\uncover<4->{\textbf{1 Define}, \textcolor{BlueLine}{\textbf{5 Deploy}}} \\
        \bottomrule
      \end{tabular}
    \end{column}
  \end{columns}
  \mlsyscredit{Symptoms: Beede et al., CHI 2020, as described in ML Lifecycle, \S3.1.2.}
  \note{[2.5 min] ENGAGE: one minute with a neighbor; each pair writes a stage number for each row. Step 2: ungradable photos are a data problem, a quality check at capture would have prompted a retake while the patient was still there, and validation on photos from real clinic cameras would have shown the gap before launch. Step 3: retakes are a workflow cost; the nurse's time is part of problem definition's operational requirements, and human-factors validation measures it. Step 4: connectivity is a deployment limit that belongs in problem definition; a design that needs the uplink for every screening fails on the day the link drops. Accept other defensible answers. WARN: if students write ``stage 3, retrain the model,'' ask what retraining would change about lighting or bandwidth. The land line for this segment: each stage hands the next an artifact, and evidence flows back, so deployment begins a phase rather than ending a project. FLEX: [CORE].}
\end{frame}

% =============================================================================
% SECTION: PROPAGATION
% =============================================================================
\mlsysnote{[0.5 min] The day-153 failure is one instance of a general principle. Name it.}
\mlsysdivider{How does one decision reach every other stage?}{ML Lifecycle, \S3.2 and \S3.4}
\section{Propagation}

\mlsysnote{[2 min] The constraint propagation principle, the key idea of the lecture. Read it aloud. A trained model, a dataset, a validation report, and a deployment package are all records of decisions. When a limit nobody stated turns up late, every artifact that silently assumed its absence must be traced, and the ones that depended on it must be redone. The size of the rework depends on which artifacts relied on the missing limit, not on a fixed multiplier for how late it was found. Then read the second line: the remedy is to state limits where they can still shape the work. WARN: students hear ``late is expensive'' and conclude that everything must be frozen early. The chapter's point is different: fix the operating envelope early (memory, latency, power, safety floors), and let data and model choices keep iterating inside it. FLEX: [CORE].}
\mlsysfocus{The Constraint Propagation Principle}{%
  A constraint discovered late forces rework in every artifact\\that relied on its absence.\\[0.5cm]
  {\normalsize\textcolor{white!75}{State a limit where it can shape the work, not where it can only reject it.}}%
}

\begin{frame}{A limit found at deployment sends work back to stage 2}
  \centering
  \mlsyssteps[width=0.92\textwidth,height=5.4cm,keepaspectratio]{propagation_step}{3}
  \mlsyscredit{ML Lifecycle, \S3.2 and Example 3.1.}
  \note{[2 min] The same loop as before. Step 1: the 500 MB limit surfaces at stage 5, deployment. Step 2: the stages whose artifacts assumed a bigger budget turn red: the preprocessing chosen in stage 2, the weights from stage 3, the validation evidence from stage 4, the deployment package from stage 5. The dashed red arrow is the rework path. Step 3: where the limit belongs, at stage 1, on day 1, where it would have shaped the choice of model instead of rejecting it. ENGAGE: ``Is stage 2 entirely red?'' Not necessarily; the next exercise asks which parts survive. FLEX: [CORE].}
\end{frame}

\begin{frame}{Constraints travel backward and forward through the stages}
  \centering
  \mlsyssteps[width=0.92\textwidth,height=5.3cm,keepaspectratio]{directions_step}{2}
  \mlsyscredit{ML Lifecycle, \S3.2.}
  \note{[2 min] Step 1, backward: the clinic's 2 Mb/s uplink is a deployment fact. If the team compresses photos to fit the link, compression leaves artifacts in the images, so the training data must contain the same artifacts, or the model sees different inputs in service than in training. That mismatch is called training-serving skew. A limit at stage 5 reached back to stage 2. Step 2, forward: an early data decision reaches later stages. The capture resolution needed to see small hemorrhages fixes the model's input size, and the input size sets the operations every inference pays on the clinic device. ENGAGE: ``Name one other backward path.'' (A latency target at deployment bounds model size at training.) Because propagation runs both ways, the stages form a network of constraints rather than a chain. LINK: Lecture 4 treats a data-quality failure that propagates downstream, a data cascade. FLEX: [CORE].}
\end{frame}

\begin{frame}{Each stage owns one axis; two stages only verify}
  \centering
  \safeimg[width=0.92\textwidth,height=5.4cm,keepaspectratio]{loop_dam.pdf}
  \mlsyscredit{ML Lifecycle, \S3.4.}
  \note{[1.5 min] The same loop, colored by the D·A·M axes from Lecture 0. Problem definition sets requirements for all three. Data collection owns the data axis, training owns the algorithm axis, deployment commits both to a machine. Evaluation and monitoring own no axis; their job is to check that the three work together, before release and then continuously. LINK: for students who remember the iron law, the chapter maps stages to its terms: data collection shapes bytes moved, training shapes the operation count and hardware efficiency, deployment shapes fixed overhead. These are starting points for diagnosis, not exclusive ownership. FLEX: [CORE].}
\end{frame}

\begin{frame}{A contract at each hand-off catches the limit on day 1}
  \centering
  \begin{minipage}{0.78\textwidth}
    \begin{mlsyscard}{mlsysaccent}
      \normalsize
      \textbf{Stage 1, problem definition}\\[0.15cm]
      \textbf{In:} business need, clinic context\\[0.05cm]
      \textbf{Out:} measurable targets, deployment tier, resource limits\\[0.05cm]
      \textbf{Invariant:} every success criterion is a number
    \end{mlsyscard}
  \end{minipage}\\[0.6cm]
  \uncover<2->{\large Deployment target: ``TBD after training'' $\rightarrow$ \textcolor{RedLine}{\textbf{blocked}}}
  \mlsyscredit{ML Lifecycle, table 2 and Example 3.1.}
  \note{[1.5 min] The chapter writes a contract for every stage boundary: what the stage needs, what it must hand over, and a quality rule the hand-off must satisfy. This is the one for problem definition. Step 2: a team that writes ``deployment target to be decided after training'' fails the audit, and data collection should not start. In production, pipeline tools encode these hand-offs as a graph of tasks that halts the moment a contract fails; the habit matters more than the tool. LINK: the same idea as an interface specification between software services. FLEX: [OPTIONAL]; table 2 of the chapter has all six contracts for students who want them.}
\end{frame}

\begin{frame}{Diagnose: a late limit. What must be redone?}
  \mlsysactivity{Diagnose}{2.5 min}
  \centering
  {\normalsize The 500 MB limit appears on day 151. Which artifacts must change?}\\[0.7cm]
  \Large
  \chip{Raw photos}{midgray}\enspace
  \chip{Labels}{midgray}\enspace
  \chip{Preprocessing}{RedLine}\\[0.35cm]
  \chip{Weights}{RedLine}\enspace
  \chip{Validation evidence}{RedLine}\enspace
  \chip{Deployment package}{RedLine}\\[0.6cm]
  \normalsize\uncover<2->{\textcolor{midgray}{\textbf{Gray}: still valid.}\quad \textcolor{RedLine}{\textbf{Red}: redo.}}
  \mlsyscredit{ML Lifecycle, Example 3.1.}
  \note{[2.5 min] ENGAGE: students write down which of the six they would redo, then compare with a neighbor. Step 2: the raw photos and the specialists' labels remain usable; changing the model does not require recollecting every image. The input preparation (resolution and preprocessing) changes for the smaller model, so the weights must be trained again, the sensitivity and on-device memory tests repeated, and a new package built. WARN: two opposite errors are common. One says ``start over'': the principle says trace dependencies, not discard everything. The other says ``just compress the 4 GB model'': compression may cut sensitivity below the floor, so validation must run again anyway. FLEX: [CORE].}
\end{frame}

\begin{frame}{Redoing four dependent artifacts takes six days}
  \centering
  \safeimg[width=0.92\textwidth,height=4.4cm,keepaspectratio]{repair_gantt.pdf}\\[0.15cm]
  {\normalsize Traced on day 1: one line in a spec. Traced on day 151: this, at best.}
  \mlsyscredit{ML Lifecycle, Example 3.1 (a supplied schedule, not a multiplier).}
  \note{[1.5 min] The chapter's worked repair: 1 day to revise preparation, 2 to train new weights, 2 to repeat validation, 1 to package. Each task waits for the artifact before it, so the chain takes 6 days if staff and compute are free. WARN: this is a supplied schedule for one dependency chain, not a rule that late discovery costs some fixed factor more; queueing, a failed validation, or a bigger model change would lengthen it. The day-153 team lost five months because nobody had traced the chain and its choices had no slack left. The land line for this segment: state a limit at problem definition, where it shapes the work, not at deployment, where it can only reject it. FLEX: [CORE].}
\end{frame}

% =============================================================================
% SECTION: ITERATION
% =============================================================================
\mlsysnote{[0.5 min] A late constraint costs time. Where does a team's time go in the first place?}
\mlsysdivider{Where do the hours go, and why does cycle time matter?}{ML Lifecycle, \S3.3}
\section{Iteration}

\begin{frame}{Predict: how many put algorithms first?}
  \mlsysactivity{Predict}{1.5 min}
  \centering
  {\normalsize A 2016 survey asked data scientists which task took most of their time.}\\[0.45cm]
  \begin{minipage}{0.55\textwidth}
    \Large
    \alt<2->{\rchoice{A}{4\% of them}}{\choice{A}{4\% of them}}\\[0.15cm]
    \choice{B}{15\%}\\[0.15cm]
    \choice{C}{30\%}\\[0.15cm]
    \choice{D}{50\%}
  \end{minipage}
  \mlsyscredit{CrowdFlower, \emph{Data Science Report}, 2016, as cited in ML Lifecycle, \S3.3.}
  \note{[1.5 min] ENGAGE: vote. Most of the room will choose C or D, because courses and papers are mostly about models. Step 2: 4 percent. The next slide shows where the rest went. FLEX: [CORE].}
\end{frame}

\begin{frame}{Data work, not model work, was the biggest time sink}
  \centering
  \safeimg[width=0.92\textwidth,height=5.0cm,keepaspectratio]{time_sinks.pdf}
  \mlsyscredit{CrowdFlower, \emph{Data Science Report}, 2016; redrawn from ML Lifecycle, figure 4.}
  \note{[1 min] Green bars are data work: cleaning and organizing (60 percent) and collecting datasets (19 percent), 79 percent together. Refining algorithms, in orange, drew 4 percent. WARN: these are shares of respondents naming one activity as their largest time sink, not shares of total hours; the chapter is careful about that. Still, in the retinopathy project every photo must be graded by several ophthalmologists before it can train anything, so collection and labeling can delay experiments more than training does. LINK: Lecture 4 is about engineering that data. FLEX: [CORE].}
\end{frame}

\begin{frame}{Six months buys 26 weekly runs or 100 hourly ones}
  \mlsysactivity{Your turn}{3 min}
  \begin{columns}[T]
    \begin{column}{0.40\textwidth}
      \begin{mlsyscard}{mlsysaccent}
        \textbf{Given: 26 weeks}\\[0.1cm]
        \textbf{Large} ensemble: 1 week per run, starts at 95\%, +0.15 per run\\[0.1cm]
        \textbf{Small} edge model: 1 hour per run, starts at 90\%, +0.1 per run\\[0.1cm]
        Ceiling: 99\%
      \end{mlsyscard}
    \end{column}
    \begin{column}{0.56\textwidth}
      \normalsize
      \uncover<2->{Large: 26 runs; $95 + 26 \times 0.15 = 98.9\%$\\[0.3cm]}
      \uncover<3->{Small: $26 \times 168 = 4{,}368$ possible runs\\
        A team can design and learn from about \textbf{100}\\[0.3cm]}
      \uncover<4->{Small: $90 + 100 \times 0.1 = 100 \rightarrow$ capped at \textbf{99\%}\\[0.3cm]}
      \uncover<5->{{\Large\textbf{\textcolor{mlsysaccent}{Same ceiling; 3.8$\times$ the ideas tested}}}}
    \end{column}
  \end{columns}
  \mlsyscredit{ML Lifecycle, Napkin Math 3.1 (assumed gains per run).}
  \note{[3 min] ENGAGE: one minute to compute the run counts before step 2. Step 2: one run per week gives 26 runs; at 0.15 points each the large model reaches 98.9 percent. Step 3: at one hour per run there are 26 times 168, 4,368 possible runs. The machine is no longer the limit; people are. The chapter counts 100 effective experiments, a few per working week, the pace at which a team can design a change, run it, and learn from it. Step 4: 90 plus 100 times 0.1 is 100, capped at the 99 percent ceiling. Step 5: both end near the ceiling, and the 0.1-point gap is an artifact of the assumed gains, not a finding. What survives is breadth: 100 against 26 experiments, 3.8 times as many ideas, so the fast team also learned which preprocessing and augmentation choices matter. WARN: this is not ``small models are more accurate.'' Speed buys search breadth, not guaranteed accuracy. FLEX: [CORE].}
\end{frame}

\begin{frame}{Cycle time sets how many ideas a team can test}
  \begin{columns}[c]
    \begin{column}{0.48\textwidth}
      \centering\safeimg[width=\textwidth,height=4.6cm,keepaspectratio]{book_iteration.pdf}
    \end{column}
    \begin{column}{0.48\textwidth}
      \normalsize
      The slow cycle is the \textbf{iteration tax}.\\[0.35cm]
      A late discovery adds to it: every rework loop is a run the team did not spend on a new idea.
    \end{column}
  \end{columns}
  \mlsyscredit{ML Lifecycle, \S3.3 (margin figure).}
  \note{[1 min] The book's figure: 100 effective experiments against 26, about 3.8 times. The chapter calls the cost of waiting on slow training runs or slow validation the iteration tax. Late discoveries add to it, because every rework loop from the last segment consumes runs that could have tested new ideas. LINK: shorter cycles come from smaller models, data selection (Lecture 9), and faster training systems (Lecture 8). The land line for this segment: cycle time sets how many ideas a team can test; a fast loop buys search breadth, not guaranteed accuracy. FLEX: [CORE].}
\end{frame}

% =============================================================================
% SECTION: DEFINE TO TRAIN
% =============================================================================
\mlsysnote{[0.5 min] Now walk the loop, one stage at a time: one picture and one failure each. The first three stages decide what the model can ever see.}
\mlsysdivider{How does a request become numbers a model can meet?}{ML Lifecycle, \S3.5 to \S3.7}
\section{Define to Train}

\begin{frame}{Stage 1: a request becomes a requirement in three layers}
  \centering
  {\large\textcolor{midgray}{``Build a model that detects diabetic retinopathy.''}}\\[0.45cm]
  \normalsize
  \begin{tabular}{@{}P{0.20\textwidth}P{0.66\textwidth}@{}}
    \leavevmode\uncover<2->{\textcolor{GreenLine}{\textbf{Statistical}}} & \leavevmode\uncover<2->{Sensitivity $>$ 90\%, specificity $>$ 80\%, in every subgroup and camera} \\[0.25cm]
    \leavevmode\uncover<3->{\textcolor{BlueLine}{\textbf{Physical}}} & \leavevmode\uncover<3->{About 500 MB, 7 to 15 W, a result in seconds, without the uplink} \\[0.25cm]
    \leavevmode\uncover<4->{\textcolor{OrangeLine}{\textbf{Operational}}} & \leavevmode\uncover<4->{Auditable decisions, regulatory clearance, the nurse's workflow} \\
  \end{tabular}
  \mlsyscredit{ML Lifecycle, \S3.5.1.}
  \note{[2 min] Step 1: the request as a product manager writes it. It hides a dozen decisions. Before any of them, the team states what improvement the system must make over the clinic's existing screening and referral process, the baseline; a learned system has to earn its maintenance cost. Step 2, statistical layer: targets over data, sensitivity above 90 percent so few sick patients are missed, specificity above 80 percent so regional specialists are not flooded with false referrals, and both must hold for every patient group and camera model, not only on average. Step 3, physical layer: the clinic device's 500 MB reservation, its 7 to 15 W power envelope, a result before the patient leaves, and operation when the uplink is down. Step 4, operational layer: staff must be able to audit why a patient was flagged, regulators must clear the device, and it must fit the nurse's day. ENGAGE: ``Which layer did the day-153 team leave out?'' (Physical.) WARN: ophthalmologists and clinic administrators must be in this meeting; engineers alone will write an accuracy target. FLEX: [CORE].}
\end{frame}

\begin{frame}{Vote: what accuracy does ``always normal'' score?}
  \mlsysactivity{Vote}{2 min}
  \centering
  {\normalsize A model answers ``normal'' for every patient. What is its accuracy?}\\[0.45cm]
  \begin{minipage}{0.55\textwidth}
    \Large
    \choice{A}{0\%}\\[0.15cm]
    \choice{B}{10\%}\\[0.15cm]
    \choice{C}{50\%}\\[0.15cm]
    \alt<2->{\rchoice{D}{90\%}}{\choice{D}{90\%}}
  \end{minipage}
  \mlsyscredit{ML Lifecycle, \S3.5.1 and \S3.8.1.}
  \note{[2 min] ENGAGE: vote, thirty seconds with a neighbor, vote again. Step 2: D, 90 percent, because it is right about every healthy patient. Students who chose A are thinking of the sick patients, which is the right instinct and the wrong metric. The next slide shows why accuracy is the wrong number for a screener. FLEX: [CORE].}
\end{frame}

\begin{frame}{Accuracy 90\%, sensitivity 0\%: screening needs two numbers}
  \begin{columns}[c]
    \begin{column}{0.60\textwidth}
      \centering\mlsyssteps[width=\textwidth,height=4.6cm,keepaspectratio]{patients_step}{2}
    \end{column}
    \begin{column}{0.38\textwidth}
      \normalsize
      \textbf{Sensitivity}\\ sick patients flagged $\div$ all sick\\[0.35cm]
      \textbf{Specificity}\\ healthy patients cleared $\div$ all healthy
    \end{column}
  \end{columns}
  \mlsyscredit{Illustrative population; floors from ML Lifecycle, \S3.5.1.}
  \note{[1.5 min] Step 1: 100 patients, 10 with disease. Step 2: the model that says ``normal'' to everyone is right 90 times, so 90 percent accuracy, and catches none of the 10, so 0 percent sensitivity. Sensitivity counts only sick patients in its denominator, so a mostly healthy population cannot inflate it. Specificity counts only healthy patients. Standard training minimizes average error, which the healthy majority dominates, so the statistical layer must say what matters: a missed case can mean blindness, a false referral costs a specialist visit. The chapter's remedies are a class-weighted loss, a sensitivity floor, and targets per subgroup. LINK: Lecture 12 covers quality metrics beyond accuracy. FLEX: [CORE].}
\end{frame}

\begin{frame}{Discuss: when the label was the bias}
  \mlsysactivity{Discuss}{3.5 min}
  \begin{columns}[c]
    \begin{column}{0.48\textwidth}
      \small
      \textbf{2014:} Amazon ranks job applicants, trained on 10 years of r\'esum\'es and past hiring decisions.\\[0.2cm]
      \textbf{Learned:} downgrade ``women's'' and graduates of two all-women's colleges.\\[0.2cm]
      \textbf{Reported 2018:} abandoned; removing gender words could not rule out other proxies.
    \end{column}
    \begin{column}{0.48\textwidth}
      \centering
      {\Large Where should this have been caught, and by whom?}\\[0.5cm]
      {\normalsize\textcolor{midgray}{Two minutes with a neighbor.}}
    \end{column}
  \end{columns}
  \mlsyscredit{ML Lifecycle, War Story 3.1, after Dastin, Reuters, 2018.}
  \note{[3.5 min] War Story 3.1. Amazon's experimental recruiting tool was trained on ten years of r\'esum\'es and the hiring decisions made about them. Because past hiring was mostly male, the model learned that male-coded language predicted success and downgraded r\'esum\'es containing ``women's'' and graduates of two all-women's colleges. Engineers removed explicit gender words, but could not guarantee the model would not find other proxies, and the tool was abandoned. ENGAGE: two minutes in pairs, then take three answers. Expected: at problem definition. The labels were past decisions, so the model learned to reproduce them; it did not learn the goal, which was good hires. Problem definition must state fairness and audit criteria and what would cause the project to be rejected, before any data is collected, and the people who own hiring policy, not only engineers, must write them. WARN: students will propose a fix at stage 3, a better model or removing a feature. That was tried and failed; the defect was in the label, which the model optimized faithfully. LINK: the same structure appears in health care, where a risk score that used past cost as its label underestimated need for one group (Obermeyer et al. 2019); Lecture 15 develops it. FLEX: [CORE].}
\end{frame}

\begin{frame}{Stage 2: clinic photos fail checks that lab photos pass}
  \begin{columns}[c]
    \begin{column}{0.48\textwidth}
      \centering\safeimg[width=\textwidth,height=4.8cm,keepaspectratio]{retinal-photo-session.jpg}
    \end{column}
    \begin{column}{0.48\textwidth}
      \centering
      {\fontsize{40}{44}\selectfont\bfseries\textcolor{RedLine}{21\%}}\\[0.2cm]
      {\normalsize of 1,838 clinic photos rejected as ungradable}\\[0.5cm]
      {\small\textcolor{midgray}{Lighting, camera upkeep, undilated pupils}}
    \end{column}
  \end{columns}
  \mlsyscredit{Photo: Staff Sgt. Jason McCasland, U.S. Air Force (public domain), Wikimedia Commons. Data: Beede et al., CHI 2020, as cited in ML Lifecycle, \S3.1.2.}
  \note{[1.5 min] Stage 2's picture: a retinal photo being taken. The result depends on the room's light, the camera's condition, whether the pupil was dilated, and the operator's training. Research datasets come from well-kept cameras under controlled conditions; clinic photos do not. That is the lab-to-field gap: the model is not wrong, the inputs have moved outside what it was trained on. Stage 2's failure, from the Thai study: 393 of 1,838 photos, 21 percent, rejected as ungradable. The chapter's fixes are a quality check at the camera, so the nurse retakes the photo while the patient is still there, and provenance on every image (camera model, settings, operator) so a problem at one site can be traced. LINK: Lecture 4 builds this pipeline. FLEX: [CORE].}
\end{frame}

\begin{frame}{Stage 3: at 224 pixels, small lesions shrink to a blur}
  \begin{columns}[c]
    \begin{column}{0.30\textwidth}
      \centering\safeimg[width=\textwidth,height=4.2cm,keepaspectratio]{retina_whole.png}\\[0.05cm]
      {\footnotesize\textcolor{midgray}{Retina with diabetic retinopathy}}
    \end{column}
    \begin{column}{0.33\textwidth}
      \centering\safeimg[width=\textwidth,height=3.6cm,keepaspectratio]{retina_full_crop.png}\\[0.05cm]
      {\footnotesize\textbf{Boxed region, as captured}}
    \end{column}
    \begin{column}{0.33\textwidth}
      \centering\uncover<2->{\safeimg[width=\textwidth,height=3.6cm,keepaspectratio]{retina_224_crop.png}\\[0.05cm]
      {\footnotesize\textbf{\textcolor{RedLine}{Same region, image shrunk to 224$\times$224}}}}
    \end{column}
  \end{columns}
  \mlsyscredit{Fundus photo: Shaofeng Hao et al. (CC BY 4.0), Wikimedia Commons. Downsampled for this slide. ML Lifecycle, \S3.7.2.}
  \note{[2 min] Stage 3's picture is what the model sees. Left: a retina with diabetic retinopathy; the white box holds small yellow deposits called exudates, a sign of leaking vessels. Middle: that region as captured. Step 2, right: the same region after the whole photo is shrunk to 224 by 224 pixels, the default input size of ImageNet models. Each lesion is now a pixel or two. Stage 3's failure: an engineer who shrinks photos to the standard size to save computation erases the evidence before training starts, and no amount of training recovers it. This is why the chapter puts ophthalmologists in model development. WARN: this public photo is already only about 1,000 pixels wide; real clinical photos have several million pixels, so the loss is larger than shown. Choosing a higher resolution then raises the operations per inference, the forward propagation from the last segment, which is why the 500 MB and latency budget must be checked at that resolution. Mention one more stage-3 failure if time allows: the Netflix Prize's winning ensemble was never deployed because its engineering cost exceeded its value. The land line for this segment: requirements are numbers in three layers, and each stage after them can destroy the evidence they need. FLEX: [CORE].}
\end{frame}

% =============================================================================
% SECTION: VALIDATE TO MONITOR
% =============================================================================
\mlsysnote{[0.5 min] The second half of the loop: the gate before release, the move into the clinic, and what changes after.}
\mlsysdivider{What must hold before release, and what changes after it?}{ML Lifecycle, \S3.8 to \S3.11}
\section{Validate to Monitor}

\begin{frame}{Stage 4: one threshold must land where both floors hold}
  \centering
  \mlsyssteps[width=0.92\textwidth,height=5.2cm,keepaspectratio]{roc_step}{3}
  \mlsyscredit{Lab point: Gulshan et al., JAMA 2016 (AUC 0.99), as cited in ML Lifecycle, \S3.7. Floors: \S3.5.1. Curve and clinic point illustrative (\S3.8).}
  \note{[2.5 min] Axes: specificity across, the share of healthy patients cleared; sensitivity up. Step 1: the two floors from stage 1, as dashed red lines; the green corner is where both hold. Step 2: a model outputs a score, and a threshold turns the score into a referral. Sweeping the threshold traces a curve: refer more patients and sensitivity rises while specificity falls. The curve here is illustrative, shaped like a model with an area under the curve (AUC) of 0.99. The blue dot is the real lab operating point from the 2016 study, 97.5 and 93.4 percent, inside the corner. Step 3: the same threshold on clinic photos, an illustrative point in the spirit of the chapter's scenario, falls below the sensitivity floor. WARN: AUC summarizes every threshold at once; the clinic runs one. A high AUC does not guarantee that any threshold meets both floors on the photos the clinic actually takes. Validation is the evidence gate: on clinic images, per subgroup (a model with 94 percent overall sensitivity can fall below 80 percent for patients with cataracts), with calibrated confidence, and with latency measured on the device. Its output is the release manifest. LINK: shadow and canary releases, Lecture 14. FLEX: [CORE].}
\end{frame}

\begin{frame}{Estimate: can a clinic upload its photos?}
  \mlsysactivity{Your turn}{2.5 min}
  \begin{columns}[T]
    \begin{column}{0.40\textwidth}
      \begin{mlsyscard}{mlsysaccent}
        \textbf{Given}\\[0.1cm]
        100 patients a day\\
        4 photos each, 5 MB per photo\\
        Uplink: 2 Mb/s, and it drops\\
        Shift: 8 hours
      \end{mlsyscard}
    \end{column}
    \begin{column}{0.56\textwidth}
      \normalsize
      \uncover<2->{$100 \times 4 \times 5\ \text{MB} = 2{,}000\ \text{MB}$ a day\\[0.3cm]}
      \uncover<3->{2 Mb/s $= 0.25$ MB/s $\Rightarrow$ 8,000 s $\approx$ \textbf{2.2 h}, 28\% of the shift\\[0.3cm]}
      \uncover<4->{One photo: $5 / 0.25 = $ \textbf{20 s}. When the link drops: never.\\[0.3cm]}
      \uncover<5->{{\large\textbf{\textcolor{mlsysaccent}{Infer in the clinic; send 10 KB summaries}}}}
    \end{column}
  \end{columns}
  \mlsyscredit{ML Lifecycle, Napkin Math 3.2.}
  \note{[2.5 min] Stage 5, deployment. ENGAGE: one minute to compute the daily volume and the upload time. Step 2: 100 patients times 4 photos times 5 MB is 2,000 MB, 2 GB a day. Step 3: 2 megabits per second is 0.25 megabytes per second, so 8,000 seconds, about 2.2 hours, 27.8 percent of an 8-hour shift. On a good day that fits. Step 4: but one photo takes 20 seconds just to upload before any remote processing, and on a day the link drops nothing uploads at all, while the nurse needs the result before the patient leaves. Step 5: the in-visit requirement, not the byte count, rules out the cloud; inference runs on the clinic device and only 10 KB summaries per patient go upstream, about 2,000 times less traffic. WARN: students learned in Lecture 2 to divide bytes by bandwidth; here the deciding fact is availability. LINK: this is the latency gate from the recall slide, and the data locality rule from Lecture 2: move the computation to the data. FLEX: [CORE].}
\end{frame}

\begin{frame}{Edge costs more up front and pays back in 2.4 years}
  \centering
  \mlsyssteps[width=0.92\textwidth,height=4.4cm,keepaspectratio]{cost_step}{3}\\[0.15cm]
  {\normalsize Cost confirms the choice that feasibility already made.}
  \mlsyscredit{ML Lifecycle, Napkin Math 3.3: 500 clinics, 50 patients a day, one graded image each; cloud fails the latency gate and is a cost benchmark only.}
  \note{[2 min] Network-wide, 500 clinics, 50 patients a day, one graded image per patient. Step 1: cloud inference at about \$0.01 per image is \$91,250 a year, plus an assumed \$45,000 a year for connectivity: about \$136,250 a year, starting from zero. Step 2: edge needs a \$500 device in each clinic, \$250,000 up front, then about \$25,000 a year of maintenance and \$9,125 a year of inference. Step 3: savings are \$102,125 a year, so the hardware pays back in \$250,000 divided by \$102,125, about 2.4 years. ENGAGE: ``What would make the lines never cross?'' (Edge running costs at or above cloud costs, or a much more expensive device.) WARN: this is the last gate, not the first. The cloud already failed the latency gate; cost here only confirms that the feasible option is also affordable. Edge also brings its own costs: models must fit the device, and updating 500 devices needs verified, atomic updates. FLEX: [CORE].}
\end{frame}

\begin{frame}{Stage 6: a new camera can lower sensitivity with no code change}
  \begin{columns}[c]
    \begin{column}{0.36\textwidth}
      \centering\safeimg[width=\textwidth,height=3.9cm,keepaspectratio]{retinal-camera.jpg}\\[0.05cm]
      {\footnotesize\textcolor{midgray}{A new camera, a new color profile}}
    \end{column}
    \begin{column}{0.62\textwidth}
      \centering\safeimg[width=\textwidth,height=4.3cm,keepaspectratio]{cadences.pdf}
    \end{column}
  \end{columns}
  \mlsyscredit{ML Lifecycle, \S3.10 and \S3.11.1 (illustrative cadences). Photo: Jason Ruck (CC BY-SA 3.0), Wikimedia Commons.}
  \note{[2 min] Stage 6's failure, from the chapter's illustrative case: six months after launch, a clinic upgrades its fundus cameras and the new photos have a different color profile. No code changed, latency and error rates stay flat, every dashboard is green, and sensitivity at that site can fall. The plot is the chapter's monitoring plan, one cadence per signal, in minutes on a log scale. A misconfigured camera shows within minutes. Proxy signals that need no labels, such as the referral rate and the model's confidence, are reviewed daily. Labeled accuracy needs confirmed diagnoses, so it arrives weekly at best. Population coverage takes months, and architecture is reviewed quarterly: about five orders of magnitude. WARN: a drift alert is evidence that inputs changed, not proof that sensitivity fell; labels decide, and labels from referred patients only cannot measure what the model missed. LINK: Lecture 14 runs this loop in production; Lecture 4 owns the drift statistics. The land line for this segment: validation tests the numbers problem definition set, deployment is screened for feasibility and then cost, and monitoring reopens the loop. FLEX: [CORE]; if short, show only the photo and say the first three sentences.}
\end{frame}

\begin{frame}{Check: true or false?}
  \mlsysactivity{Vote}{2 min}
  \normalsize
  \renewcommand{\arraystretch}{1.4}
  \begin{tabular}{@{}P{0.66\textwidth}P{0.27\textwidth}@{}}
    \toprule
    \textbf{Claim} & \textbf{Verdict} \\
    \midrule
    Passing evaluation means the system is ready to deploy. & \leavevmode\uncover<2->{\textcolor{RedLine}{\textbf{False}}: 21\% ungradable} \\
    The deployment target can be chosen after training. & \leavevmode\uncover<3->{\textcolor{RedLine}{\textbf{False}}: 4 GB vs 500 MB} \\
    Skipping validation stages speeds delivery. & \leavevmode\uncover<4->{\textcolor{RedLine}{\textbf{False}}: risk moves to patients} \\
    \bottomrule
  \end{tabular}
  \mlsyscredit{ML Lifecycle, \S3.12, Fallacies and Pitfalls.}
  \note{[2 min] Three of the chapter's fallacies and pitfalls. ENGAGE: thumbs up or down for each, then reveal. Step 2: false; a model with specialist-level accuracy on curated images rejected 21 percent of clinic photos, and evaluation never tested serving, uploads, or workflow. Step 3: false; the day-153 failure, a 4 GB model against a 500 MB reservation. Step 4: false; a skipped gate does not remove the risk it checked, it moves it into production, where a failure means containment, diagnosis, and validation again under pressure. A gate may be removed only when its evidence is supplied another way. Mention the fourth if time allows: more data is not always the best investment, because examples from regions already well covered add little, while targeted collection can fix a failing subgroup. FLEX: [CORE].}
\end{frame}

% =============================================================================
% SYNTHESIS AND CLOSE
% =============================================================================
\begin{frame}{Back to the clinics: every failure had a stage and a fix}
  \begin{columns}[c]
    \begin{column}{0.36\textwidth}
      \centering\safeimg[width=\textwidth,height=4.4cm,keepaspectratio]{loop_step3.pdf}
    \end{column}
    \begin{column}{0.62\textwidth}
      \normalsize
      \uncover<1->{\textcolor{GreenLine}{\textbf{Ungradable photos}}: check quality at the camera; validate on clinic photos.\\[0.3cm]}
      \uncover<2->{\textcolor{BlueLine}{\textbf{Stalled uploads}}: state connectivity at problem definition; infer in the clinic.\\[0.3cm]}
      \uncover<3->{\textbf{Google's own change}: specialists review ungradable photos with the patient's records.}
    \end{column}
  \end{columns}
  \mlsyscredit{Google, ``Healthcare AI systems that put people at the center,'' 2020; Beede et al., CHI 2020; ML Lifecycle, \S3.1.2 and \S3.6.}
  \note{[2 min] Return to the opening question with today's tools. Step 1: the ungradable photos belong to data collection, where a quality check at capture prompts a retake while the patient is still present, and to validation, which must use photos from the cameras and rooms the system will meet. Step 2: the stalled uploads belong to problem definition, where the clinic's connectivity is a stated limit, and to deployment, which then runs inference in the clinic. Step 3: Google's team changed the workflow, not the model: rather than referring every patient with an ungradable photo, eye specialists now review those photos alongside medical records. ENGAGE: ``Which of these fixes required retraining?'' (None.) Close with the thesis: the system around the model decided the outcome, and seeing the whole loop first is how a team finds these failures on day 1 instead of in the clinic. FLEX: [CORE].}
\end{frame}

\begin{frame}{Takeaways}
  \small
  \begin{itemize}\setlength{\itemsep}{0.12cm}
    \item The lifecycle is a loop of six stages; monitoring sends evidence back, so deployment starts a phase.
    \item A late constraint forces rework in every artifact that relied on its absence: trace the dependencies.
    \item Cycle time sets how many ideas a team can test; speed buys breadth, not accuracy.
    \item Requirements are numbers in three layers: statistical, physical, operational.
    \item Accuracy is not readiness: validation tests the floors on field data, per subgroup, on the device.
    \item Screen deployment for feasibility first; cost only confirms.
  \end{itemize}
  \note{[1.5 min] One line per idea. Ask students to circle the one they could not yet explain to a friend and reread that section first. FLEX: [CORE].}
\end{frame}

\begin{frame}{Open questions}
  \normalsize
  \begin{itemize}\setlength{\itemsep}{0.4cm}
    \item Who should own the deployment limit: the ML team, the hardware team, or the clinic?
    \item When labels arrive only for patients the model referred, how can a team measure what it misses?
    \item Should an approved medical model be allowed to change, and how often?
  \end{itemize}
  \note{[1 min] No single answer. Assign one as a short written reflection. The second has a partial answer in the chapter: a sample of images, sent off-hours and labeled centrally, supplies outcome evidence that referrals alone cannot.}
\end{frame}

\begin{frame}{Read and explore}
  \small
  \begin{itemize}\setlength{\itemsep}{0.2cm}
    \item \textbf{Read:} ML Lifecycle, \S3.1 to \S3.5 and \S3.8 to \S3.10; redo Napkin Math 3.1 to 3.3 by hand.
    \item Beede et al., ``A Human-Centered Evaluation of a Deep Learning System Deployed in Clinics for the Detection of Diabetic Retinopathy,'' CHI 2020.
    \item Gulshan et al., ``Development and Validation of a Deep Learning Algorithm for Detection of Diabetic Retinopathy in Retinal Fundus Photographs,'' JAMA 2016.
    \item Amershi et al., ``Software Engineering for Machine Learning: A Case Study,'' ICSE-SEIP 2019.
    \item Sculley et al., ``Hidden Technical Debt in Machine Learning Systems,'' NeurIPS 2015.
  \end{itemize}
  \note{[1 min] The chapter first, with its self-check questions. Each outside reading backs one part of today: the clinic study that is the hook (Beede), the lab result and how sensitivity and specificity are reported (Gulshan), the ML workflow as practiced at Microsoft, which the chapter's lifecycle definition cites (Amershi), and why decisions made in isolation accumulate cost (Sculley).}
\end{frame}

\begin{frame}[c]
  \centering
  {\small\textcolor{lightgray}{NEXT LECTURE \textbar{} DATA ENGINEERING}}\\[0.5cm]
  {\LARGE\bfseries Why does training data function\\like source code?}
  \note{[0.5 min] Today the loop; next, its first material. The retinopathy screener needs about 100,000 graded photos, quality checks at the camera, and ingestion that survives a fragile uplink. Lecture 4 treats data as the program and engineers it accordingly. Ask students to read Data Engineering before then.}
\end{frame}

\end{document}
